\section{Résumé du cas clinique}

\subsection{Motif}
La patiente se présente aux urgences, le $30$ décembre $2025$, accompagnée de la police pour une demande d’internement à la suite d'une tentative de suicide.

\subsection{Contexte social}
La patiente a $60$ ans, elle a un fils de $27$ ans. Elle s’est séparée du père lorsque celle-ci était encore enceinte. Depuis madame vit seule avec son fils.

Au niveau professionnel, elle est pensionnée depuis le $1$ janvier $2026$. Elle a fait carrière en tant que secrétaire de direction.

La patiente rapporte avoir été victime de violence dans son enfance. 
Ensuite, son ancien compagnon était violent avec elle. De plus, l’enfant est issu d’un viol de celui-ci. Elle a toutefois, décidé de se séparer de cet homme tout en gardant l’enfant et en l’élevant en tant que mère célibataire. Actuellement elle subit des violences physiques et verbales de la part de son fils rendant la situation intenable au domicile. Celui-ci a été diagnostiqué dès son plus jeune âge de trouble de l’attention et d’hyperactivité ainsi que du trouble du spectre autistique. En effet, il présente un comportement violant à la moindre contrariété. Celui-ci a un suivi psychiatrique mais n’a pas de situation professionnelle ni logement.

\subsection{Antécédents psychiatriques}
\begin{itemize}[label=\textbullet]
    \item \underLine{Depuis l’âge de $18$ ans} : dépression récurrente.
    \item  \underLine{De $2015$ à $2021$ } : suivi psychiatrique mais aucun suivi ambulatoire depuis.
    \item \underLine{ $2021$} : dépendance aux benzodiazépines.
    \item \underLine{ $2021$ et juillet $2025$} : $2$ hospitalisations pour dépression.
\end{itemize}

\subsection{Antécédents somatiques}
\begin{itemize}[label=\textbullet]
    \item Actuellement en rémission d’un mélanome.
    \item Thyroïdite d’Hashimoto  
    \item Fibromyalgie
    \item Vessie hypertonique
    \item Diarrhées chroniques idiopathiques
\end{itemize}

\subsection{Allergies}
\begin{itemize}[label=\textbullet]
    \item Moxifloxacine 
\end{itemize}
 
\subsection{Antécédents familiaux}
\begin{itemize}[label=\textbullet]
    \item Son père affecté d’un cancer prostatique au stade métastatique. 
\end{itemize}

\subsection{Anamnèse}  
Suite à une tentative d’étranglement de la part de son fils, la patiente aurait fait une tentative de suicide. La patiente dit avoir pris $5-6$ comprimés de lormétazépam $2$mg, $6-7$comprimés de lorazépam $1$mg, $2$ comprimés d'alprazolam retard dont la dose est inconnue. Elle aurait été retrouvée inconsciente à l’arrivée de la Police qui avait été prévenue par les voisins. Elle souhaite inscrire son fils auprès d’associations qui pourraient l’aider à trouver un logement et un travail mais les délais d’attente sont longs. La patiente ne voit pas de solutions possibles à court terme. Malgré la thymie triste et les ruminations suicidaires qui persistent, la patiente se projette dans l’avenir et a des projets. 

\subsection{Examen clinique}
La patiente est anxieuse avec une thymie triste et labile.
Elle est cohérente. Une hypoprosexie est observée. 
Son discours est volubile et circonlocutoire pouvant témoigner d’une grande charge anxieuse associé à des difficultés d’attention et une accélération psychique. 
La patient présente des insomnies de maintien. 
Sur le plan psycho-affectif, on note une intense ambivalence vis-à-vis du fils. 
La patiente oscille entre un désir de protection maternelle, teinté de culpabilité, et un besoin de mise à distance vitale face aux hétéro-agressions subies. Cette dynamique relationnelle complexe alimente un sentiment d'impasse existentielle, majorant le risque suicidaire. 
Une hyper-émotivité est observée chez la patiente ce qui se perçoit comme une personne qui s’implique beaucoup avec les personnes qui sont bienveillantes avec elle. 

\subsection{Examens complémentaires}
\begin{itemize}[label=\textbullet]
    \item Biologie : sans particularité mis à part une hypercholestérolémie à $243$mg/dl. (Ionogramme et fonction rénale (>$60$ mL/min) sont corrects).
    \item ECG : rythme sinusal, FC $80$bpm
\end{itemize}
Voir détails des examens complémentaires en \cref{annexe:examens_complementaires} en \cpageref{annexe:examens_complementaires}.

\subsection{Diagnostic retenu}
 
Épisode dépressif majeur avec caractéristiques mixtes\footnote{\url{https://www.infodrog.ch/files/content/refbases/DSM-5_Manuel-diagnostique-et-statistique-des-troubles-mentaux.pdf\#page=286} \parencite{dsm5_online}}  dans un contexte familial relationnel difficile.

\subsection{Traitements réalisés et prise en charge} 
\begin{itemize}[label=\textbullet]
    \item Paroxétine (Seroxat)  $20$mg $1$cp/j
    \item Zopiclone (Imovane) $7,5$mg $1$cp/j
    \item Olanzapine (Zyprexa) $5$mg $1$cp/j 
    \item Lorazépam (Temesta) $2,5$mg $1$cp/j
    \item Lévothyroxine sodique (T4, 100 µg) et liothyronine sodique (T3, 20 µg) (Euthyral) $75$mg $1$cp/j 
    \item Lopéramide (Imodium) $2$mg $2$cp/j
    \item Fésotérodine (Toviaz) $4$mg $1$cp/j 
\end{itemize}
 
Ce traitement pharmacologique est associé à une psychothérapie et l’hospitalisation sera suivi d’une cure psychothérapeutique stationnaire. 
Une assistante sociale s’intéresse également au cas afin de tenter d’apporter des solutions dans la sphère sociale de la patiente. 
 
Lors de l’hospitalisation, un changement au niveau du traitement a eu lieu. Le valproate de sodium $500$mg a été remplacé par l’olanzapine $5$mg principalement à cause d’une chute de cheveux marquée lors de l’augmentation de la posologie du valproate. De l’olanzapine est mis à la place avec effet positif thymorégulateur et nette atténuation des fluctuations émotionnelles au cours d’une même journée. 
 
L’abandon de la quétiapine $20$mg a également eu lieu car elle a présenté des éruptions cutanées. 

\subsection{Évolution de l’hospitalisation }
Amélioration thymique et amélioration de son sommeil. 
La patiente a des idées noires passagères toutefois dans l’évolution de son hospitalisation, celle-ci demande de l’aide auprès des soignants. 
La patiente apprend a posé ses limites auprès de son entourage. 

\subsection{Ressenti de la patiente}
La patiente semble être demandeuse de traitement et a une volonté d’aller mieux. La patiente voit son hospitalisation comme un moyen se protéger de son fils et comme une ressource sociale pour l’aider dans les démarches administratives et trouver des solutions.
 
Concernant son traitement psychiatrique, la patiente associe cela à un même besoin qu’un traitement somatique. Elle a toujours considéré sa santé mentale au même titre que sa santé physique. Il est donc normal pour elle de prendre un traitement quotidiennement. 
La patiente se décrit comme une personne «hypersensible» pharmacologiquement parlant, en d’autres termes, avec le temps, elle a eu l’occasion de tester de nombreuses molécules différentes et selon elle, elle ressentirait les effets secondaires plus fréquemment que la population générale et avec des doses plus faibles. Au cours de sa prise en charge avec son psychiatre extérieur, elle dit n’avoir jamais trouvé l’équilibre entre bénéfices et effets secondaires, y compris avec le valproate. Depuis l’instauration de l’olanzapine, elle se sent mieux. 